% ===========================================================================
%  Qubit & Toffoli Reduction Techniques: A Comprehensive Primer
%  LaTeX source — generated from optimization_techniques_primer.md
% ===========================================================================
\documentclass[11pt,a4paper]{article}

% ---- Encoding & fonts ----
\usepackage[T1]{fontenc}
\usepackage[utf8]{inputenc}
\usepackage{lmodern}
\usepackage{microtype}

% ---- Page layout ----
\usepackage[top=2.5cm,bottom=2.5cm,left=2.8cm,right=2.8cm]{geometry}
\usepackage{fancyhdr}
\pagestyle{fancy}
\fancyhf{}
\rhead{\small\textit{Qubit \& Toffoli Reduction Techniques}}
\lhead{\small\textit{ecdsa.fail Benchmark Primer}}
\cfoot{\thepage}
\renewcommand{\headrulewidth}{0.4pt}

% ---- Math ----
\usepackage{amsmath,amssymb}

% ---- Colors ----
\usepackage[dvipsnames,svgnames]{xcolor}
\definecolor{codebackground}{gray}{0.95}
\definecolor{codeborder}{gray}{0.75}
\definecolor{tableheadbg}{RGB}{230,240,255}
\definecolor{accent}{RGB}{0,70,140}
\definecolor{warnbg}{RGB}{255,248,230}

% ---- Hyperlinks ----
\usepackage[colorlinks=true,linkcolor=accent,urlcolor=accent,citecolor=accent,
            pdftitle={Qubit \& Toffoli Reduction Techniques: A Comprehensive Primer},
            pdfauthor={ecdsa.fail research notes},
            bookmarks=true]{hyperref}

% ---- Code listings ----
\usepackage{listings}
\usepackage{fancyvrb}

\lstset{
  basicstyle=\ttfamily\footnotesize,
  breaklines=true,
  breakatwhitespace=false,
  frame=single,
  rulecolor=\color{codeborder},
  backgroundcolor=\color{codebackground},
  xleftmargin=6pt,
  xrightmargin=6pt,
  framexleftmargin=4pt,
  framexrightmargin=4pt,
  aboveskip=8pt,
  belowskip=6pt,
  showstringspaces=false,
  captionpos=b,
  columns=flexible,
  keepspaces=true,
  tabsize=2,
}

\lstdefinelanguage{Rust}{
  keywords={fn,let,mut,if,else,match,while,for,in,return,pub,use,mod,struct,enum,
             impl,trait,type,where,self,Self,true,false,Some,None,Ok,Err,usize,u8,u16,u32},
  keywordstyle=\color{NavyBlue}\bfseries,
  comment=[l]{//},
  commentstyle=\color{gray}\itshape,
  stringstyle=\color{ForestGreen},
  morestring=[b]",
}

% ---- Tables ----
\usepackage{booktabs}
\usepackage{array}
\usepackage{longtable}
\usepackage{tabularx}
\usepackage{multirow}
\renewcommand{\arraystretch}{1.25}

% ---- Colored boxes ----
\usepackage[most]{tcolorbox}

\tcbuselibrary{breakable,skins}

\tcbset{
  definitionbox/.style={
    breakable,
    enhanced,
    colback=blue!4,
    colframe=accent,
    fonttitle=\bfseries,
    top=4pt, bottom=4pt, left=6pt, right=6pt,
  },
  noteboxstyle/.style={
    breakable,
    enhanced,
    colback=warnbg,
    colframe=orange!60!red,
    fonttitle=\bfseries,
    top=4pt, bottom=4pt, left=6pt, right=6pt,
  },
}

% ---- Enumeration ----
\usepackage{enumitem}
\setlist[itemize]{leftmargin=1.5em,itemsep=1pt,topsep=3pt}
\setlist[enumerate]{leftmargin=1.5em,itemsep=1pt,topsep=3pt}

% ---- Section styling ----
\usepackage{titlesec}
\titleformat{\section}{\large\bfseries\color{accent}}{\thesection.}{0.7em}{}[\color{accent}\titlerule]
\titleformat{\subsection}{\normalsize\bfseries}{\thesubsection.}{0.5em}{}
\titleformat{\subsubsection}{\normalsize\itshape\bfseries}{}{0em}{}
% Number subsections as letters within their section (6.A, 6.B, ... 8.A, ...) to
% match the in-text references and avoid redundant "8.1 8.A" double labels.
\renewcommand{\thesubsection}{\thesection.\Alph{subsection}}

% ---- Utilities ----
\usepackage{xspace}
\usepackage{parskip}
\setlength{\parindent}{0pt}
\setlength{\parskip}{5pt}
\setlength{\emergencystretch}{3em}   % allow extra inter-word stretch to avoid overfull boxes

% ---- Custom commands ----
\newcommand{\ket}[1]{$|#1\rangle$\xspace}
\newcommand{\score}{\textsc{score}\xspace}
\newcommand{\sota}{\textsc{sota}\xspace}
\newcommand{\cls}{\texttt{cls}\xspace}
\newcommand{\pha}{\texttt{pha}\xspace}
\newcommand{\anc}{\texttt{anc}\xspace}
\newcommand{\tof}{Toffoli\xspace}
\newcommand{\avgT}{\texttt{avg\_executed\_Toffoli}\xspace}

\setcounter{tocdepth}{2}

% ===========================================================================
\begin{document}
% ===========================================================================

\begin{titlepage}
\centering
\vspace*{2cm}
{\LARGE\bfseries\color{accent} Qubit \& Toffoli Reduction Techniques\par}
\vspace{0.4cm}
{\Large A Comprehensive Primer\par}
\vspace{1.2cm}
{\large\itshape ecdsa.fail quantum circuit benchmark\par}
\vspace{0.5cm}
{\normalsize Based on research notes from the trailmix\_ludicrous optimization track\par}
\vspace{2.5cm}
\begin{tcolorbox}[definitionbox, title=Audience note]
An undergrad who has learned the basics of quantum computing --- superposition, entanglement,
Toffoli/CNOT/Hadamard gates, quantum circuits, and reversible computation.
This document explains every optimization technique used in the ecdsa.fail challenge
from circuit-level first principles, with concrete examples from actual research.
\end{tcolorbox}
\vfill
{\small As of 2026-06-27\par}
\end{titlepage}

\tableofcontents
\newpage

% ===========================================================================
\section{What We Are Optimizing}
\label{sec:score}
% ===========================================================================

\subsection{The score}

\begin{tcolorbox}[definitionbox, title=Score objective (lower is better)]
\[
  \text{score} = \text{peak\_qubits} \;\times\; \text{avg\_executed\_Toffoli}
\]
\end{tcolorbox}

This is the ecdsa.fail benchmark's cost function for a quantum circuit that computes an affine
elliptic-curve point addition on secp256k1 (Bitcoin's curve). The two factors penalize two
different resources:

\begin{itemize}
  \item \textbf{Peak qubits}: the maximum number of qubits simultaneously in use at any instant.
    This is the ``width'' of the circuit, not the total ever allocated. On real hardware, idle
    qubits still occupy physical space on the quantum processor.
  \item \textbf{Average executed \tof}: the Toffoli (CCX) gate count, \emph{weighted by how
    often each gate fires}, averaged over the 9024 test inputs the challenge uses. Not just the
    gate count, but gates~$\times$~firing\_rate.
\end{itemize}

\subsection{Why Toffoli gates specifically?}

On a \textbf{fault-tolerant} quantum computer (the kind needed to break Bitcoin-scale
cryptography), gates fall into two classes:

\begin{itemize}
  \item \textbf{Clifford gates} (CNOT, Hadamard $H$, Phase $S$, Pauli $X/Y/Z$): implemented
    efficiently using the surface error-correcting code at very low physical-qubit overhead.
  \item \textbf{Non-Clifford gates} (\tof / T-gate): require \textbf{magic states} --- specially
    prepared ancilla qubits that take a long time and many physical qubits to distill via
    ``magic state distillation.'' Each \tof gate consumes one distilled magic state.
\end{itemize}

The rough conversion: $1\ \text{Toffoli} = \text{exactly}\ 7\ T\text{-gates}$ (Gosset et al.\
2013, tight lower bound).

So in this regime, \textbf{\tof gates are like currency}: each one costs approximately the same
fixed price (one magic state). The peak qubit count is like ``desk space'' --- idle qubits still
occupy factory floor. The product $\text{peak\_qubits} \times \text{Toffoli}$ captures the
spacetime volume of running the circuit.

\begin{tcolorbox}[noteboxstyle, title=Common metric confusion]
This challenge uses \avgT, NOT:
\begin{itemize}
  \item T-count ($= 7 \times$ Toffoli, a different metric)
  \item T-depth (Selinger 2013 --- a depth trick with 4 ancilla per Toffoli; irrelevant here)
\end{itemize}
The distinction matters: average-executed rewards conditional execution (replay), emitted
T-count does not.
\end{tcolorbox}

\subsection{The current SOTA}

The best known result as of 2026-06-27 (BitWonka, commit \texttt{d44cad3},
trailmix\_ludicrous family):
\[
  1152\ \text{qubits} \times 1{,}364{,}230\ \text{avg Toffoli} = 1{,}571{,}592{,}960\ \text{score}
\]

\subsection{Two separate competitions, two separate objectives}
\label{sec:twocomps}

This is important context for the whole document. There are \textbf{two distinct optimization
tracks}, and a technique that is optimal for one can be irrelevant for the other:

\begin{enumerate}
  \item \textbf{The Q$\times$T score competition} (the main subject of this primer): minimize the
    \emph{product} $\texttt{peak\_qubits} \times \texttt{avg\_executed\_Toffoli}$. Here both
    resources matter, and every move is judged by the exchange rate (\S\ref{sec:exchange}). The SOTA
    above (1152q $\times$ 1.36M) lives in this track.
  \item \textbf{The low-qubit competition} (a separate objective): minimize \textbf{peak qubits
    alone}. Toffoli count is \emph{unconstrained} --- you may spend as many Toffoli gates as you
    like. The only thing that counts is how few qubits the circuit is simultaneously live on.
\end{enumerate}

These are genuinely different games. Under the Q$\times$T objective you would never spend 8$\times$
more Toffoli to save 10\% of qubits; under the low-qubit objective that trade is \emph{free}, because
Toffoli simply does not appear in the scoring. Most of this primer optimizes track 1, but the
deepest qubit-floor technique --- \textbf{EEA register sharing (\S\ref{sec:regshare})} --- belongs
to track 2, and reaches qubit counts (851q $\to$ 829q) far below anything competitive on the product
score. Keep the two objectives separate in your head: many ``is this worth it?'' questions only have
an answer once you know \emph{which} competition you are playing.

% ===========================================================================
\section{The Reversibility Constraint}
\label{sec:reversibility}
% ===========================================================================

\subsection{You cannot delete information}

Quantum circuits must be reversible. This is not a convention --- it is required by quantum
mechanics. Every gate in a quantum circuit must have an inverse (you can always ``run it
backwards''). This means:

\begin{itemize}
  \item \textbf{You cannot overwrite a qubit.} If qubit $A$ holds value $x$ and you want to
    compute $f(x)$ into $A$, you cannot just replace it. You'd lose $x$ forever, breaking
    reversibility.
  \item \textbf{You cannot erase a qubit.} If a qubit is entangled with the rest of the system
    (as intermediate results always are), setting it to \ket{0} forcibly would violate unitarity.
\end{itemize}

The only way to ``free'' a qubit is to \textbf{uncompute} it --- run the computation that created
its value in reverse, returning it exactly to \ket{0}. Once it's \ket{0}, it can be freed.

\subsection{Consequence: intermediate values cost qubits until uncomputed}

Every time your circuit computes an intermediate result (a carry bit, a partial product, a
comparison flag), it must allocate a fresh qubit. That qubit remains \textbf{live} --- occupying
a slot in the peak count --- from the moment it's computed until the moment it's uncomputed.

\textbf{Example}: An $n$-bit addition $a + b$ computes $n$ carry bits, one per bit position. Each
carry qubit is live from when it's computed until the carry chain is reversed. If you compute the
addition at step 100 and uncompute it at step 200, those $n$ carry qubits are live across the
entire steps 100--200, contributing to the peak at any instant in that range.

This is why almost every qubit-reduction technique boils down to:
\textbf{make intermediate values die sooner} (so they aren't live at the peak) \textbf{or don't
materialize them at all}.

\subsection{Reversibility also doubles Toffoli counts}

In a classical circuit, you compute $X$, use $X$, then discard $X$. In a reversible circuit, you
compute $X$ (costs \tof), use $X$, then \emph{uncompute} $X$ (costs the same \tof again). So
every logical AND operation normally costs $2\times$ --- once forward, once backward.

This is the central pain point that Measurement-Based Uncomputation (\S\ref{sec:mbu}) solves.

% ===========================================================================
\section{The Challenge Circuit}
\label{sec:circuit}
% ===========================================================================

\subsection{What it computes}

The circuit computes one \textbf{affine elliptic curve point addition}: given quantum register
$P = (P_x, P_y)$ and classical constant $Q = (Q_x, Q_y)$, compute $P + Q = R$ in place, where
the arithmetic is modulo the secp256k1 prime $p = 2^{256} - 2^{32} - 977$.

The mathematical operations, in order:
\begin{lstlisting}
dx = Px - Qx              (mod p subtraction)
dy = Py - Qy
lambda = dy / dx          (modular inverse + multiply: the expensive GCD step)
new_Px = lambda^2 - Px - Qx
new_Py = lambda*(Px - new_Px) - Py
\end{lstlisting}

\subsection{Why modular inverse is the bottleneck}

Computing $dy/dx \bmod p$ requires computing $dx^{-1} \bmod p$, the modular inverse of $dx$.
There is no simple formula --- it requires running the \textbf{Extended Euclidean Algorithm (EEA)}
or equivalent, which takes $\approx530$ steps of binary GCD (divstep) operations. Each step involves:
\begin{itemize}
  \item A comparison (swap decision: is $|u| > |v|$?)
  \item A conditional swap of two 256-bit registers
  \item A conditional subtraction of two 256-bit registers
\end{itemize}
This accounts for roughly 90\% of the circuit's \tof count.

\subsection{The key design choice: Q is classical}

The second point $Q$ is a \textbf{fixed constant} (the secp256k1 generator point, known at circuit
compile time). The current SOTA keeps $Q$ as classical bit-patterns (\texttt{BitId}s) rather than
quantum registers, eliminating \textbf{512 qubits} from the peak --- the single largest qubit
reduction in the circuit's history.

% ===========================================================================
\section{Measuring Progress}
\label{sec:measuring}
% ===========================================================================

\subsection{The live-qubit-vs-time profile}

To reduce the peak, you first need to know \textbf{where} the peak occurs. Plot how many qubits
are live at each moment in the circuit timeline:

\begin{lstlisting}
qubits  |
live    |         ##########
        |    ####################
        |  #########################
        |################################
        +--+----+--+------+--+----+----> time
                    ^
                    peak = 1152 here (the "binder")
\end{lstlisting}

The highest point is the \textbf{peak}, and the phase executing at that height is the
\textbf{binder}. Only operations alive at the binder instant matter for qubit reduction.

\textbf{Key tool}: \texttt{TRACE\_PEAK=1 TRACE\_PHASE\_ACTIVE=1} in the build flags. This
instruments every \texttt{alloc\_qubit}/\texttt{free\_qubit} call and prints the name of the
phase executing when the maximum is hit.

\subsection{The plateau problem}

After shaving the single tallest binder, you often discover the peak hasn't moved --- because
several \emph{independent} phases all reach the same height. This is a \textbf{co-binder plateau}:

\begin{lstlisting}
qubits  |
live    |    #####   #####     ####
        |  #########################
        |  #########################  <- plateau at 1152 (Phase A, B, C all tie here)
        +----+--+----+---+--+---+---> time
\end{lstlisting}

The reported peak doesn't drop until \textbf{every} plateau member is individually pushed below the
current height. This changes the strategy fundamentally: you must work in ``plateau mode'' ---
find a local qubit reduction for each co-binder, then combine them all.

\subsection{Emitted vs average-executed Toffoli}

\begin{itemize}
  \item \textbf{Emitted}: the raw count of CCX gates in the circuit, regardless of whether they fire.
  \item \textbf{Average-executed}: each CCX gate's contribution weighted by $\Pr[\text{control fires}]$,
    averaged over all 9024 test inputs. A CCX that only fires 5\% of the time contributes only 0.05.
\end{itemize}

This distinction is critical because \textbf{conditional replay} (\S\ref{sec:replay}) can
dramatically reduce average-executed without changing emitted at all.

% ===========================================================================
\section{Qubit Reduction: The Core Loop}
\label{sec:qubitloop}
% ===========================================================================

Before diving into individual techniques, understand the general loop:

\begin{enumerate}
  \item \textbf{Measure} --- get the live-qubit-vs-time trace; find the peak height and the binder phase.
  \item \textbf{Inventory} --- list every qubit alive at the binder instant.
  \item \textbf{Classify} --- label each qubit by category (table below).
  \item \textbf{Brainstorm} --- for each non-essential qubit at the binder, propose a shave.
  \item \textbf{Cost} --- compute the \tof overhead vs qubits saved. Compare to the break-even exchange rate (\S\ref{sec:exchange}).
  \item \textbf{Re-measure} --- apply the cheapest safe shave. The peak moves to a new binder. Repeat.
\end{enumerate}

\subsection{Qubit classification at the binder}

\begin{center}
\begin{tabular}{p{3.2cm}p{5.5cm}p{3.5cm}}
\toprule
\textbf{Class} & \textbf{What it holds} & \textbf{Shrinkable?} \\
\midrule
Essential data & Inputs/outputs the binder reads or writes & Rarely \\
Scratch & Temporary arithmetic results (carries, partial products) & Often \\
Transcript / log & GCD branch decisions recorded for reversibility & Often --- compress the log \\
Passenger & Allocated before the peak, consumed after, \emph{idle at binder} & Yes --- relocate its endpoints \\
Alias-candidate & Restorable duplicate of a value already live elsewhere & Yes --- borrow as dirty scratch \\
Truncatable & Bits provably zero or irrelevant to the final output & Yes \\
\bottomrule
\end{tabular}
\end{center}

\subsection{Binary vs graded failure modes}

A critical distinction before applying any technique:

\begin{itemize}
  \item \textbf{Graded-failure levers} (carry truncation, comparator narrowing): correctness
    degrades \emph{continuously} with aggressiveness. A small fraction of inputs go wrong. You can
    iterate on the parameter and search for a nonce where all 9024 test inputs avoid the failures.

  \item \textbf{Binary-failure levers} (register live-range holes, transcript compression,
    structural reuse): correctness is all-or-nothing. Either the circuit is value-exact on all
    inputs, or it's wrong on essentially \emph{all} inputs. There is no graded near-miss ladder
    to climb. You must get the implementation mathematically correct before searching for nonces.
\end{itemize}

\begin{tcolorbox}[noteboxstyle, title=The trap]
Applying the ``tweak-and-triage'' loop to a binary-failure lever.
A structural width cut that's not value-exact will produce all-dirty garbage on every input, not
near-misses.
\end{tcolorbox}

% ===========================================================================
\section{Qubit Reduction Techniques}
\label{sec:qubittech}
% ===========================================================================

\subsection{Live-Range Holes: Uncompute Early, Recompute Later}
\label{sec:liverange}

\textbf{The headline move.} If a qubit holds a value $V$ that the peak instant does NOT need (it
was computed before the peak and will be consumed after), free it before the peak and rebuild it
after.

\begin{lstlisting}
Timeline view:
Before:  ...[V computed]=====[V held live across peak]=====[V consumed]...
After:   ...[V computed][uncompute V]  (peak happens)  [recompute V][V consumed]...
                                        ^^^^^^^^^^^
                                  V is absent at peak
\end{lstlisting}

\textbf{Mechanics}: Emit the reverse of the gates that built $V$ (returning its qubits to \ket{0}
and freeing them), run the peak-owning phase with that freed headroom, then re-emit the forward
gates to restore $V$ for its later consumer.

\textbf{Cost}: For a value that cost $C_0$ Toffoli to compute:
\begin{itemize}
  \item Save $k$ qubits at the peak
  \item Spend $2C_0$ extra Toffoli (one uncompute + one recompute)
  \item The trade wins if $2C_0 < k \times (\text{avgT} / \text{peak\_q})$ (the break-even rate)
\end{itemize}

\textbf{Real examples}:
\begin{itemize}
  \item \textbf{1211$\to$1193}: A GCD-apply scratch block was uncomputed during a shift sub-phase
    that didn't use it, then reacquired after. Saved 18 qubits.
  \item \textbf{1166$\to$1165} (\texttt{TLM\_PARK\_ODD\_U0}): \texttt{u[0]} is provably 1 (GCD
    odd-u invariant). Apply $X \to$ \ket{0} to zero it, \texttt{loan\_zero\_qubit} to return the
    lane, then \texttt{restore\_known\_one} (apply $X$ again) to restore. Near-zero \tof cost,
    $-1$ qubit.
\end{itemize}

\textbf{The keep-alive dual}: Sometimes the uncompute+recompute round-trip creates a wider
transient intermediate than just holding the value. In that case, \emph{keep} the value live
across the middle operation. Choose whichever is narrower: holding footprint vs recompute peak
transient.

\subsection{Passenger Relocation}

A \textbf{passenger} is a qubit that was allocated early and consumed late, but does nothing at
the peak --- it just sits idle.

\textbf{The cheaper fix}: Don't uncompute and recompute. Just \textbf{allocate later} (right
before first use) and \textbf{free earlier} (right after last use). If both endpoints fall
outside the peak window, the passenger vanishes from the binder count at zero \tof cost.

\textbf{Check passengers first} before resorting to live-range holes --- relocation is free.

\subsection{Range-Bound a Register to Drop a Guard Bit}

Some registers need an extra ``guard bit'' (sign or overflow indicator) only because an
intermediate could transiently go negative or overflow. If you can prove the value always stays
in a non-negative range by reordering operations, delete the guard bit entirely.

\textbf{Example} (benhuang Lever B): The fold quotient register had bit 33 (a sign guard) because
Solinas reduction terms could make the partial sum dip below zero. Reorder to apply positives
first. Now the sum is always in $[0, 2^{33})$. Allocate 33 bits instead of 34.
\textbf{Net: $-1$ qubit at zero \tof cost.} This produced the 1221$\to$1220 drop.

\textbf{General rule}: In $+A - B$ sequences, applying positives first and subtracting last
keeps partial sums non-negative, eliminating sign-guard bits.

\subsection{Truncation: Keep Only the Bits That Matter}

Carry chains, high words of partial sums, and intermediate registers often hold more bits than
actually affect the final result.

\textbf{Per-step scheduling}: In a 530-step loop (the GCD), a single constant truncation width
must be set safe for the worst iteration. A per-step schedule $\text{width}(\text{step})$ lets
each iteration be exactly as tight as it can tolerate --- saving qubits only at the steps that
are actually at the peak.

\textbf{Two truncation flavors}:
\begin{itemize}
  \item \textbf{Structural truncation (density-neutral)}: prove from circuit analysis that the
    dropped bits are always zero on all inputs.
  \item \textbf{Empirical truncation (island-exact)}: verify the dropped bits are never 1 on the
    9024 challenge inputs specifically. Requires re-hunting a nonce after each change.
\end{itemize}

\subsection{Free-and-Recompute a Cheap Value (The 1153$\to$1152 Breakthrough)}

When a dead-qubit scan finds \textbf{no} unused qubits at the peak, don't conclude
``the floor is structural.'' Ask instead:

\begin{quote}
  \textit{Is any held value a cheap function of the other values already live at the peak?}
\end{quote}

If yes, \textbf{free} that qubit, let the peak phase use its slot, then \textbf{recompute} the
value from the live qubits afterward for negligible cost.

\textbf{The cy0 trick} (\texttt{d44cad3}, BitWonka): The qubit \texttt{cy0} holds the carry-in
at bit 0. No dead qubits found by scan. But because the secp256k1 prime constant $f[0] = 1$:
\[
  \texttt{cy0} = \texttt{control\_bit}\ \text{AND}\ (\text{NOT}\ \texttt{final\_a0})
\]
Both \texttt{control\_bit} and \texttt{final\_a0} are already live at the peak. So:
\begin{enumerate}
  \item Uncompute \texttt{cy0} using 2 CNOT-equivalent gates ($\approx 0$ Toffoli)
  \item Return the freed lane via \texttt{loan\_zero\_qubit}
  \item After the peak phase: recompute \texttt{cy0} from \texttt{control\_bit} and
    \texttt{final\_a0} (2 gates)
\end{enumerate}
Applied 40 times (once per binding fold iteration): $40 \times 4$ CNOTs $= 0$ Toffoli.
Saves 1 qubit. The 1153$\to$1152 drop.

\textbf{Generalized lesson}: After a dead-qubit scan fails, scan again for ``whose value is a
CHEAP FUNCTION of other live qubits?'' The second question catches what the first misses.

\subsection{Dynamic-Width Registers}

In a GCD algorithm, operands naturally shrink as computation proceeds. After $k$ steps of binary
GCD, the operands $u$ and $v$ have lost at most $k$ bits from their high ends. A naive circuit
allocates full 256-bit registers for all 530 steps.

\textbf{Dynamic-width registers}: At each step $i$, determine the actual live bit-width and only
allocate that many qubits. Free the high bits as they become known-zero.

\textbf{Ablation results} (from the Shrunken-PZ q980 track):
\begin{lstlisting}
Source-level dynamic PZ (universal schedule):     1027 qubits
+ thin support-tuned schedule (per-step widths):   990 qubits  (-37q)
+ counter/rotation metadata trimming:              983 qubits   (-7q)
+ quotient cap:                                    980 qubits   (-3q)
\end{lstlisting}

\textbf{Thin schedule (island-exact)}: An even tighter version samples actual GCD widths over many
random inputs and builds per-step bounds tight to that distribution. This is island-exact ---
the thin schedule may be too narrow for inputs outside the sampled set. Requires nonce hunting.

\subsection{Known-Constant Teardown Before the Peak}

When the GCD terminates, certain registers are in known, predetermined states: $A = 0$, $B = 1$,
$\texttt{ca} = p$, $q = 0$. These are \textbf{classical constants} --- not quantum information.

\textbf{What to do}: XOR the known constant \emph{out of} the register (zeroing it) before the
peak computation, then XOR it \emph{back in} afterward.

\textbf{Example}: Register \texttt{ca} holds $p = 2^{256} - 2^{32} - 977$ at termination. Before
the lambda multiply: XOR $p$ out of \texttt{ca}, returning it to \ket{0}. Free the register. Run
the lambda multiply with the freed headroom. After the lambda multiply: XOR $p$ back in for
uncomputation.

\textbf{Toffoli cost}: XOR (CNOT) operations are Clifford gates --- free.
The teardown and recreation cost 0 Toffoli. This is a pure qubit win.

\subsection{Transcript / Log Compression}
\label{sec:transcript}

\textbf{The GCD must record its decisions.} At each of the 530 divstep iterations, the algorithm
makes data-dependent decisions: did we swap? did we subtract? These decisions are recorded in a
\textbf{transcript} (also called ``dialog tape'') because they're needed to run the GCD backward
during uncomputation. With one qubit per decision, the raw transcript is $\approx772$ bits long.

\textbf{Compression techniques}:

\textbf{1. Base-5 codec} (trailmix\_ludicrous): Each 3-step window of $(\texttt{subtracted},
\texttt{swap}, s_2)$ has 8 possible bit patterns, but only 5 are reachable. Encode each 3-step
window in $\lceil \log_2(5^3) \rceil = 7$ bits instead of 9 raw bits.
\begin{lstlisting}
Raw:        1 + 3*257 = 772 qubits
Compressed: 1 Step0 (2b) + 85 Triples (7b each) + 1 Pair tail (5b) + t1 flag (1b) = 603 qubits
\end{lstlisting}

\textbf{2. K5 exact 11-bit codec}: Enumerate all reachable 5-step transcript windows
exhaustively. If there are exactly $2^{11} = 2048$ reachable patterns, encode each in 11 bits.
Information-theoretically optimal. Implemented via a borrowed-ancilla 17-CCX reversible codec.
Converts 1203q$\to$1193q.

\textbf{3. Three-step tail codec} (32 symbols): The last few GCD steps are heavily constrained.
Enumerate the exact reachable set (32 symbols), encode in 5 bits. The 1185q$\to$1170q drop.

\textbf{4. Streaming decompression}: Instead of decompressing the entire transcript before use,
decompress one window at a time, use it, then immediately clear that window. Only one window is
``open'' at a time, reducing peak overhead from the codec.

\textbf{The hard constraint}: The codec must be a \textbf{bijection} --- a perfect 1-to-1 map
between reachable patterns and codes. Every codec must include a self-test proving bijectivity over
the full reachable support.

\subsection{Plateau Decomposition: Lower Every Co-Binder Together}

When multiple independent phases all tie at the same peak height, you're in plateau mode.
Lowering one phase alone doesn't change the global peak.

\textbf{The protocol}:
\begin{enumerate}
  \item List ALL phases tied at the current peak (not just the first one). Build a ``co-binder ledger.''
  \item Find a \emph{local} qubit reduction for each binder independently.
  \item Verify each local reduction doesn't accidentally raise a neighboring phase.
  \item Combine all compatible local reductions in one package.
  \item Re-measure the combined package.
\end{enumerate}

\textbf{Real example --- the q1170 plateau} (9-way co-bind, all tied):
\begin{center}
\begin{tabular}{p{4.5cm}p{5.5cm}p{3cm}}
\toprule
\textbf{Phase} & \textbf{Local fix} & \textbf{Effect} \\
\midrule
Solinas square & smaller square segment notch & square peak 1170$\to$1169 \\
apply final ripple & low-qubit final ripple variant & ripple 1170$\to$1162 \\
apply double/reverse halve & fold carry parking + per-step map & 1170$\to$1169 \\
special overflow/underflow folds & fold parking + per-step map & 1170$\to$1169 \\
non-final apply ripple & source-as-carry suffix for one lane & 1170$\to$1169 \\
\bottomrule
\end{tabular}
\end{center}
Each was independently developed, then combined. Only when all binders were locally reduced
did the global peak drop to 1169.

\subsection{Dynamic Headroom Clamp}

Define a \textbf{circuit-wide qubit ceiling} and have every arithmetic primitive clamp its
carry/chunk width to stay under it:

\begin{lstlisting}[language=Rust]
fn target_qubit_headroom(circ: &Circuit) -> usize {
    TLM_TARGET_Q - circ.active_qubits()
}

fn varchunk_adder(circ: &mut Circuit, bits: usize) -> ... {
    let k = min(scheduled_k, target_qubit_headroom(circ));
    // ... use k as the carry chunk width ...
}
\end{lstlisting}

\textbf{Effect}: No single call can exceed \texttt{TLM\_TARGET\_Q} total qubits.
\textbf{Lower \texttt{TLM\_TARGET\_Q} by 1 $\to$ peak drops by 1}, as long as the narrower
adder is still correct.

\textbf{The tradeoff}: Tighter ceilings force narrower carries $\to$ more carry-chunk splits $\to$
more Toffoli. The break-even exchange rate governs when to tighten.

\textbf{Crucial nuance --- the clamp is only a \emph{ceiling}, not a lane-freeing lever.} Lowering
\texttt{TLM\_TARGET\_Q} declares an aspiration; it does nothing unless there are \emph{actually
freeable lanes} for the narrower arithmetic to fit into --- ``a clamp with no freed lanes just
panics or fails to fit.'' So a clamp drop must be \textbf{paired with companion levers that genuinely
eliminate live qubits} at the peak. The \textbf{1156 $\to$ 1153} drop is the textbook example --- it
stacked three coordinated levers:

\begin{enumerate}
  \item \textbf{The clamp} itself: \texttt{TLM\_TARGET\_Q} 1156 $\to$ 1152 (the binding peak lands
    at 1153, one above the aspiration --- the irreducible co-resident set).
  \item \textbf{Zero-cin fold} (\texttt{TLM\_FOLD\_CHUNK\_ZERO\_CIN}): the fold-chunk loop's
    persistent carry-in ancilla \texttt{cin0} is \emph{no longer allocated at all} --- the first
    chunk's carry-in is provably 0, so a new ``zero-cin boundary-erase'' path
    (\lstinline{fold_boundary_erase_zero_direct}) does the carry cleanup without it. This is the
    actual freed lane (a live-range elimination, cf.\ \S\ref{sec:lazycarry}) that the clamp reclaims.
  \item \textbf{FFG vent-count cap} (\texttt{TLM\_FFG\_MAX\_G = 47}): caps the FFG half-product
    fold's clean-vent count $g$ so fewer transient vent qubits co-reside at the peak
    ($\texttt{capped\_g} = \min(\texttt{scheduled\_g}, \texttt{cap})$).
\end{enumerate}

\textbf{And the qubit drop alone is not enough to \emph{win}.} Pushing the peak 1156 $\to$ 1153
makes the clamped adders run on tighter chunks $\to$ more carry work $\to$ avgT went \emph{up} by
$+11{,}369$ (to 1{,}392{,}603), i.e.\ $\sim$3{,}790 Toffoli/qubit, far above break-even. That
``clean q1153 base'' was \textbf{rejected} until Toffoli levers (the cinc cascade replacement
\S\ref{sec:cinc} + iterated dead-CCX \S\ref{sec:deadccx}) were stacked back on top to carry it over
the line --- a perfect illustration of the product-race rule (\S\ref{sec:exchange}): a qubit drop is
only real on the score once enough Toffoli is recovered to pay for it.

\subsection{Lazy Carry Liveness}
\label{sec:lazycarry}

A specific instance of live-range shortening for carry-in qubits. The carry-in qubit was being
allocated before a chunk loop and held live across all $N$ iterations, even though it was only
needed at the loop boundaries.

\textbf{Fix} (\texttt{TLM\_FOLD\_CHUNK\_LAZY\_CIN0}):
\begin{lstlisting}
Before: carry_in allocated -> live across all N chunk iterations -> used only at boundaries
After:  carry_in allocated inside boundary_erase() -> used -> freed inside boundary_erase()
\end{lstlisting}
The carry-in qubit's liveness is reduced to just the boundary operation.
Peak saved: 1 qubit per chunk at 0 Toffoli cost.

\subsection{In-Place Decode: Hold Blocks Compressed}

For persistent register blocks that live across the peak, encode them compressed and decode
\emph{in place} only at the single use site.

\textbf{Before}: Full-width raw K2-pair block (6 lanes) held live across the whole peak.
\textbf{After}: 5 compressed cells + one zero lane stored at rest; decode in place only at the
use site (allocate the 6 lanes transiently, use them, re-encode to 5 cells). During the rest of
the circuit, only 5 cells are allocated. This is \texttt{DIALOG\_GCD\_K2\_APPLY\_INPLACE\_RAW\_BLOCK=1},
the 1218$\to$1215 qubit drop.

\subsection{Borrowed-Carry Fusion}
\label{sec:borrowcarry}

Standard carry/borrow vectors allocate fresh qubits. But if an adjacent scratch register is
currently known \ket{0}, borrow its qubits as the carry vector instead:

\begin{lstlisting}[language=Rust]
borrowed_const_fold_carries(circ, need, borrowed_carries);
// Uses borrowed[..need] as carry vector; no fresh allocation
// At exit: borrowed restored to |0>
\end{lstlisting}

\textbf{Net effect}: No fresh qubits allocated for the carry prefix. The peak doesn't rise.

\textbf{Hard requirement}: The borrowed qubits must be genuinely \ket{0} at the borrow point and
restored to \ket{0} before the borrower reads them again as data. Getting this wrong causes
silent data corruption. (\texttt{DIALOG\_GCD\_SPECIAL\_FOLD\_BORROW\_CARRIES=1}.)

\subsection{Passenger Ghosting via HMR}

\textbf{Standard passenger relocation}: Free before peak; recompute after. Requires the value
to be recomputable from live data (costs Toffoli).

\textbf{HMR ghost} (more aggressive): Measure the passenger qubit in the Hadamard ($X$-basis).
The qubit is freed. The measurement result (a classical bit $m$) is recorded. Later, reconstruct
the value algebraically and apply a classically-conditioned $Z$ correction to fix the phase.

\textbf{Why this works}: An $X$-basis measurement of $|\psi\rangle$ projects it into $|+\rangle$
or $|-\rangle$. The measurement outcome $m$ tells you which. If $m = 1$ (got $|-\rangle$), a
phase $(-1)$ was applied to the computational-basis components --- the CZ correction cancels this
exactly. After the correction, the measured qubit is factored out and can be freed without
information loss.

\textbf{In the q980 EC point addition}:
\begin{itemize}
  \item $dy$ (256-bit) computed before the GCD inversion. After GCD, $\lambda$ and $dx$ are live.
  \item The identity $dy = \lambda \cdot dx$ holds by construction.
  \item HMR-ghost $dy$ before the GCD peak (measure it, record $m$).
  \item At the end: reconstruct $dy = \lambda \cdot dx$, apply phase correction via $m$.
  \item Freed 256 qubits across the peak at the cost of one modular multiply for reconstruction.
\end{itemize}

\subsection{EEA Register Sharing --- The Low-Qubit Frontier}
\label{sec:regshare}

\textbf{The deepest qubit-reduction idea, and the headline technique for the low-qubit competition}
(\S\ref{sec:twocomps}, track 2 --- minimize peak qubits, Toffoli unconstrained). This is what drives
the lowest qubit counts ever recorded on the challenge: the 948q $\to$ 851q $\to$ 829q descent.
Because that competition does not score Toffoli at all, register sharing is free to spend gates
lavishly in exchange for lanes --- exactly the trade it makes.

\textbf{The naive EEA register budget.} The extended Euclidean algorithm carries four full-width
256-bit registers through all $\sim$530 steps:
\begin{itemize}
  \item $A$, $B$ --- the two operands (the shrinking $u$, $v$)
  \item $CA$, $CB$ --- their two B\'ezout cofactors (the running $a$, $b$ with $a\cdot dx + b\cdot p = \gcd$)
\end{itemize}
Naively that is $4 \times 256 = 1024$ qubits just for the inversion core, before any scratch.

\textbf{The sharing observation} (Proos--Zalka 2003, the origin of ``inversion dominates point-add
space''): the operands and their cofactors are \emph{complementary in size}. As the algorithm runs,
operand $A$ loses high bits (it shrinks toward 1) while its cofactor $CA$ grows from zero. At every
step the sum $\mathrm{bitlen}(A) + \mathrm{bitlen}(CA)$ is bounded by roughly one word. So $A$ and
$CA$ can be \textbf{packed into a single physical 256-bit register}: $A$ in the high part, $CA$ in
the low part, with a small \emph{bit-length tracker} recording where the boundary sits. As $A$
shrinks and $CA$ grows, the boundary simply slides --- but the total never exceeds one word.

\begin{lstlisting}
Naive:   [====== A ======][== CA (mostly zero) ==]   two separate 256-bit words
         [====== B ======][== CB (mostly zero) ==]

Shared:  [== A ==|== CA ==]   ONE word; boundary slides right as A shrinks, CA grows
         [== B ==|== CB ==]   ONE word; a 9-bit length tracker remembers the split
\end{lstlisting}

Four full registers collapse to \textbf{two ``work words'' + a handful of small length/control
trackers}. This is the structural source of Luo et al.\ 2026's compact exact reversible inversion at
$3n + 4\lfloor \log_2 n \rfloor$ qubits (1333q for $n=256$ as a standalone inversion). In the
Shrunken-PZ track, packing $A\Vert CA$ and $B\Vert CB$ plus dirty-catalytic borrowing and
gate-hosting drove the inversion core down to
\lstinline{REGISTER_SHARED_PAPER_INVERSION_QUBITS = 3*256 + 52 = 820} qubits, giving an
\textbf{851-qubit peak} for the whole point-add --- a 97-qubit drop below the prior 948q record, and
the 851$\to$829q race then squeezed out another 22.

\textbf{Where the Toffoli goes (and why that's fine here).} Packing is not free in gates: a shared
word has a \emph{data-dependent boundary} --- ``where does $A$ end and $CA$ begin?'' depends on
values computed at runtime. So every arithmetic step must first \textbf{locate and realign that
boundary with a full-width comparator}, and because the circuit is reversible, that comparator is
recomputed \emph{before and after} each hosted gate:

\begin{lstlisting}
Cost to realize ONE useful Toffoli on a packed word:
   toggle_gate ; ccx ; toggle_gate
   where toggle_gate ~= a full-width comparator ~= 12 x width ~= 3,075 Toffoli at width 256
\end{lstlisting}

Across 530 divsteps this lands the 851q route at $\sim$464.5M average Toffoli (vs $\sim$54.8M for
the unshared 948q route). \textbf{Under the low-qubit objective this is a non-issue} --- Toffoli does
not enter the score, so spending $\sim$8$\times$ more of it to buy 97 qubits is precisely the
intended trade. (For contrast, under the \emph{Q$\times$T product} objective the very same packing
would be $\sim$250$\times$ over break-even, so this technique belongs squarely to the low-qubit
track, not the score track. The two objectives have genuinely different optimal constructions --- see
\S\ref{sec:exchange} for the exchange-rate math behind that statement.)

\begin{tcolorbox}[noteboxstyle, title={A note on rigor --- is 851q a \emph{real} qubit count?}]
Only partly. It is a \textbf{lower-bound witness}, not a proven width, and the gap is worth
understanding precisely (the source spells it out):
\begin{itemize}
  \item \textbf{The deepest figure is a \emph{classical} oracle, not a circuit.}
    \lstinline{register_shared_eea.rs} computes on \texttt{U256}/\texttt{U512} integers --- \textbf{not
    a single qubit or gate in it}. Its own header: \emph{``an analysis oracle, not a reversible
    circuit implementation. In particular, passing this oracle does not establish a challenge-valid
    qubit count.''} It only checks the packing invariant $l_t + 1 + l_q + \mathrm{bitlen}(r) \le n+3$
    at all 1{,}479 steps --- i.e.\ that operand and cofactor \emph{fit in one word}, proving the
    packing is feasible in principle, classically.
  \item \textbf{The low count is bought by gate-hosting whose zero-claims are \emph{sampled, not
    proven}} (see \S\ref{sec:gatehost}). \lstinline{q945_local_hosts.rs}: \emph{``does not certify
    reachable-state zero claims. Q945 acceptance remains blocked until the Q946 hardening
    prerequisites $\ldots$ are integrated''}; \lstinline{q944_gate_host_lifecycle.rs} checks the host
    composition only \emph{``over every basis state at small widths,''} then extrapolates to 256 bits.
  \item \textbf{The submission self-reports rejection.}
    \texttt{Q949\_ROBUST\_ENVELOPE\_FRESH\_VALIDITY\_CLAIMED = false} (with an \texttt{assert!}
    enforcing it), and every artifact carries \texttt{"proof\_status":"reject"}.
\end{itemize}
So passing the $0/0/0$ run check certifies the \textbf{output values are correct on the hunted
island}; it does \textbf{not} certify that 851 qubits suffices for \emph{all} inputs. The lowest
count backed by a fully-emitted, island-validated circuit \emph{without} the extra uncertified-hosting
assumption is the \textbf{948q} route. Treat 851q/829q as ``this width looks reachable, modulo an
uncertified zero-lane assumption.'' The reusable, sound idea underneath is the structural one ---
operand/cofactor packing.
\end{tcolorbox}

\textbf{The lesson.} Register sharing is \emph{the} route below $\sim$948 qubits, and the
Proos--Zalka / Luo-2026 mechanism is the structural source worth knowing. It is the right tool when
your objective is pure qubit count. The broader point generalizes both ways: \textbf{know which
competition you are playing before pricing a lever} --- a move that is a clear win for the low-qubit
objective (trade unlimited Toffoli for lanes) can be a clear loss for the Q$\times$T objective, and
vice versa.

\subsection{Gate-Hosting: Borrowing an Idle Lane Instead of Allocating}
\label{sec:gatehost}

\textbf{The technique that powers the 948 $\to$ 851 drop (\S\ref{sec:regshare}) --- and the clearest
place to see how a qubit lever can quietly become island-overfit at the level of a \emph{single
borrowed qubit}.}

\textbf{The idea.} Every reversible gadget needs \textbf{scratch ancilla} (a multi-controlled-X
needs a lane to accumulate the AND of its controls; a carry chain needs carry bits; a comparator
needs a borrow bit). Allocating a fresh $|0\rangle$ for each is what drives the peak up. But at the
instant you need scratch, some lane \emph{elsewhere} is often idle --- the high bits of a shrunk
number, a register between its last read and next write, an idle metadata lane.
\textbf{Gate-hosting borrows that idle lane as scratch, then restores it before its owner needs it
again} --- so the ancilla costs no new peak qubit. The scratch is ``hosted'' inside a register that
belongs to something else (the codebase has explicit per-step tables naming which \texttt{A/B/CA/CB/Q}
lane+bit hosts each gate). This generalizes \S\ref{sec:borrowcarry} (borrowed-carry fusion) to
\emph{any} gate's scratch.

\textbf{The hot-desking intuition.} Qubits are desks in a crowded office. Allocating a fresh ancilla
$=$ renting a new desk (the office gets more crowded). Gate-hosting $=$ doing your hour of work at a
colleague's empty desk while they are at lunch, then clearing it spotless before they return --- no
new desk, the peak never rose.

\textbf{Two flavors:}
\begin{itemize}
  \item \textbf{Clean hosting} --- the borrowed lane is known $|0\rangle$. Use it directly, restore
    it to $|0\rangle$.
  \item \textbf{Dirty hosting} --- the borrowed lane holds unknown live data $\psi$. You can still
    borrow it via a trick that disturbs $\psi$ and then exactly un-disturbs it (the Barenco
    ``surrounded'' pattern):
\begin{lstlisting}
ccx(c0, c1, psi)        <- scribble onto the borrowed lane
ccx(psi, c2, target)    <- do the real work using it
ccx(c0, c1, psi)        <- un-scribble: psi restored to its original value
ccx(psi, c2, target)
\end{lstlisting}
    The two extra gates restore $\psi$ perfectly. Cost: 4 Toffoli instead of 3 --- a textbook
    \textbf{$-1$ qubit, $+1$ Toffoli} trade (great when the objective does not score Toffoli).
\end{itemize}

\textbf{It is a compile-time bet, not a runtime scan.} A quantum circuit is static --- you cannot
peek at a qubit mid-run to check ``is this $|0\rangle$?'' (measuring would collapse it). So nothing
searches for zeros at runtime. The hosting choices are \textbf{frozen into the gate list at build
time}: the compiler asserts \emph{``at step 437, lane 12 of register B will be $|0\rangle$, so host
the scratch there,''} and the circuit runs that exact plan on every input. The whole question is
whether that compile-time claim was \emph{true}.

\textbf{Proven-zero vs sampled-zero (value-exact vs island-exact, one qubit at a time).} The
zero-claims come in two grades:
\begin{itemize}
  \item \textbf{Proven zero (sound, density-neutral):} the lane is zero for a \emph{structural}
    reason --- e.g.\ the high bits of an operand the GCD has provably shrunk, or a known-constant
    register. Safe for all inputs.
  \item \textbf{Sampled zero (island-exact, risky):} the claim ``this lane is $|0\rangle$ here'' is
    established by an \textbf{empirical census} --- run many sample inputs, observe the lane was
    always zero, \emph{assume} it always is. This is the grade the aggressive 851q borrows rely on.
    ``The desk was empty every time I checked'' is not ``the desk is always empty.''
\end{itemize}

\textbf{How it breaks.} If an input arrives (one not in the sample) that makes a sampled-zero host
lane \emph{non-zero} at borrow time, the hosted gate scribbles onto live data and the circuit
silently returns a wrong answer --- reversibility offers no safety net. The reason the shipped
circuit still scores $0/0/0$ is that the \textbf{nonce/test island is hunted to dodge exactly those
inputs}: the 9024 graded inputs are chosen so none trips a bad borrow. So it never breaks on its own
test set --- because the test set was fitted to the circuit, not because the circuit is universally
correct.

\textbf{Takeaway.} Gate-hosting is a perfectly sound, standard peak-reduction technique \textbf{when
the host lane's state is proven}; it becomes an island-overfit assumption the moment the state is
only \emph{sampled}. The same value-exact vs island-exact split from \S\ref{sec:density} applies
here at the granularity of a single borrowed $|0\rangle$.

% ===========================================================================
\section{Toffoli Reduction: The Core Loop}
\label{sec:tofloop}
% ===========================================================================

\begin{enumerate}
  \item \textbf{Measure} --- get the \tof count; find the \textbf{hot-spot}. Note emitted vs average-executed.
  \item \textbf{Classify} --- identify what \emph{kind} of Toffoli these are (table below).
  \item \textbf{Pick the matching move} --- the kind determines the technique.
  \item \textbf{Cost it} --- depth-neutral? qubit cost? worth it at the exchange rate?
  \item \textbf{Re-measure} --- the hot-spot migrates; repeat.
\end{enumerate}

\textbf{Classifying Toffoli by kind:}
\begin{center}
\begin{tabular}{p{6.5cm}p{6cm}}
\toprule
\textbf{Kind of Toffoli at the hot-spot} & \textbf{The matching move} \\
\midrule
Carry-uncompute (reverse carry chain returning scratch to \ket{0}) & Vent via MBU (\S\ref{sec:mbu}) \\
Self-uncomputing scratch (comparators built then cleanly reversed) & Conditional replay (\S\ref{sec:replay}) \\
Data-controlled arithmetic (\texttt{cadd}, fold-sub on live data) & Fusion / avoid materialization \\
Majority gadgets (3-CCX majority in carry logic) & Ancilla-free majority (\S\ref{sec:majority}) \\
Redundant final cleanup & Skip it (\S\ref{sec:skipclean}) \\
Full intermediate products / wide rows & Avoid materialization (\S\ref{sec:witness}) \\
Adjacent algebraic stages (add/negate/subtract chains) & Algebraic fusion (\S\ref{sec:fusion}) \\
\bottomrule
\end{tabular}
\end{center}

% ===========================================================================
\section{Toffoli Reduction Techniques}
\label{sec:toftech}
% ===========================================================================

\subsection{Measurement-Based Uncomputation (MBU)}
\label{sec:mbu}

\textbf{The \tof lifecycle problem}: An AND gate costs 4 $T$-gates (1 Toffoli) to compute. The
uncomputation normally costs another 4 $T$-gates. Computing + uncomputing one AND $= 2$ Toffoli.

\textbf{MBU} (Jones 2013, applied to adder carry chains by Gidney 2018):
\begin{lstlisting}
1. Forward: ccx(a, b, out)       -> 4 T-gates (1 Toffoli)
2. Use 'out' in downstream logic
3. Uncompute: measure 'out' in X-basis (Hadamard, then Z-basis measure)
   Result is classical bit m in {0, 1}
   - If m = 0: done, free 'out'  -> 0 Toffoli
   - If m = 1: apply CZ(a, b)    -> 0 Toffoli (CZ is a Clifford gate)
Total uncompute cost: 0 Toffoli
\end{lstlisting}

\textbf{Why does this work?} When you compute \texttt{out = a AND b} into a clean \ket{0}
ancilla, the system is in state $|a, b, a\cdot b\rangle$. Measuring \texttt{out} in the
$X$-basis either succeeds cleanly ($m=0$) or applies a phase kick $(-1)$ to the
$|a=1, b=1\rangle$ component ($m=1$). The CZ gate corrects exactly that phase. After
correction, \texttt{out} is disentangled from $(a,b)$ and can be freed.

The key insight: \textbf{measurement collapses entanglement for free}. Coherent uncomputation
must undo entanglement gate-by-gate; measurement does it in one shot.

\textbf{The tradeoff}: The \texttt{out} qubit must remain live from when it's computed to
when it's measured (typically between the forward and reverse carry chains). This costs $+1$
peak qubit during that window.

\textbf{Practical effect on adders}: An $n$-bit ripple-carry adder normally costs $\approx 2n$
Toffoli ($n$ forward, $n$ reverse). With MBU on every carry-uncompute:
\[
  \text{Total adder cost} = n\ \text{Toffoli (forward only)} + 0\ \text{(measured uncompute)} = n\ \text{Toffoli}
\]
This halves the adder cost --- the most important constant-factor improvement in the SOTA.

\subsection{The Vent Dial}

MBU gives a precise, tunable dial. In a hybrid carry adder with $k$ vents:
\begin{align*}
  \text{Toffoli cost} &= 3n - 2 - k \\
  \text{Peak qubits}  &= \text{baseline} + k \quad \text{(held only during forward$\leftrightarrow$reverse window)}
\end{align*}
So each vent is exactly: \textbf{$-1$ Toffoli, $+1$ temporary peak qubit}.

\textbf{When to vent}:
\begin{itemize}
  \item \textbf{Off-peak phases}: vent aggressively. The $+1$ qubit is off-peak, no score impact.
    Pure savings.
  \item \textbf{On-peak phases}: don't vent, or only if the \tof savings outweigh the qubit cost
    at the current exchange rate.
\end{itemize}

In \texttt{trailmix\_ludicrous}, the schedule bakes vent budgets per call: each call's vent count
is set to exactly \texttt{TLM\_TARGET\_Q} $-$ \texttt{active\_qubits}, spending every spare qubit
below the ceiling on \tof savings. Calls at the ceiling get 0 vents.

\begin{tcolorbox}[noteboxstyle, title=Gidney 2025 streaming vented adder (open)]
arXiv:2507.23079 uses 2 clean $+$ $(n-2)$ dirty ancilla for $\approx 3n$ Toffoli. Partially
implemented in \texttt{venting.rs} but \textbf{currently leaks phase} --- dirty ancilla phase
correction logic is incomplete. Expected $\approx n/4$ \tof savings per adder once fixed.
\end{tcolorbox}

\subsection{Conditional / Measured Replay (Average-Executed Only)}
\label{sec:replay}

\textbf{The idea}: If a comparator flag is \ket{0} on most inputs, measure the flag instead of
keeping it coherent. Then replay the dependent computation conditionally on the measurement outcome.

\begin{lstlisting}
1. Compute flag F into scratch ancilla (costs C Toffoli)
2. Measure F in X-basis -> classical bit m
3. If m = 0: dependent block is a no-op on this shot  -> 0 Toffoli from the block
4. If m = 1: replay the block classically conditioned on m  -> full cost
Average cost = C + Pr[F=1] * (block cost)
\end{lstlisting}

When $F$ fires on fraction $f$ of inputs: average-executed drops from $(\text{block cost})$ to
$f \times (\text{block cost})$. If $f = 0.05$, you pay only 5\% on average.

\textbf{The hard constraint}: This ONLY works on \textbf{self-uncomputing scratch} --- a comparator
you compute, read under the condition, and cleanly reverse. It does NOT work on data-controlled
operations or predicates reused downstream.

\subsection{Algebraic Fusion: Collapse Add/Negate Chains}
\label{sec:fusion}

Adjacent modular-arithmetic stages that are algebraically equivalent to a simpler single
operation can be collapsed to eliminate intermediate carry chains.

\textbf{Example} (\texttt{DIALOG\_FUSE\_C\_FORM}):
\begin{lstlisting}
Old:
  tx += 2*Qx    (mod add)
  tx = -tx      (mod negate)
  tx -= Qx      (mod sub)
  tx = -tx      (mod negate)

Algebraically (two negations cancel, constants combine):
  tx += 3*Qx    (one mod add)
\end{lstlisting}

\textbf{Example} (\texttt{DIALOG\_FUSE\_X\_RESTORE}):
\begin{lstlisting}
Old:
  tx = -tx      (mod negate)
  tx += Qx      (mod add)

Algebraically:
  tx = Qx - tx  (one operation)
\end{lstlisting}

Each removed negate/add/subtract eliminates a carry chain, a modular reduction fold, and a
cleanup phase --- saving significant Toffoli.

\subsection{Witness-First / Deferred-Output Dataflow}
\label{sec:witness}

Sometimes a cancellation only needs a \emph{witness relation}, not the final output value.
Computing the output early and then cleaning it is wasteful when the witness is cheaper.

\textbf{The 4352cfb Toffoli cut} ($-10.8$M Toffoli at fixed 980 qubits):

Key identity (from curve equation): $\texttt{new\_dy} = \lambda \cdot \texttt{new\_dx}$.
So to cancel $\lambda$, we do not need to compute \texttt{new\_y} first. Instead:
\begin{enumerate}
  \item Erase \texttt{dy} using: $dy = \lambda \cdot dx$ (subtract $\lambda \cdot dx$ from
    $dy \to$ zero)
  \item Reuse the zeroed \texttt{dy} register as scratch
  \item Compute \texttt{new\_x}
  \item Compute $\texttt{new\_dx} = ox - \texttt{new\_x}$
  \item Compute $\texttt{new\_dy} = \lambda \cdot \texttt{new\_dx}$ (the witness, not the output)
  \item Cancel $\lambda$ using $\texttt{new\_dy}/\texttt{new\_dx}$
  \item Materialize $\texttt{new\_y} = \texttt{new\_dy} - oy$ (only NOW compute the output)
\end{enumerate}

\textbf{Ablation}:
\begin{lstlisting}
zero-dy disabled, defer-y disabled:   110,227,695 emitted ops
defer-y materialization only:         105,573,887 emitted ops  (-4.6M)
both zero-dy + defer-y:                92,105,748 emitted ops  (-18.1M total)
\end{lstlisting}

\textbf{General principle}: Before materializing any output, ask: ``does the next operation NEED
the output, or just a function of it that's cheaper to compute directly?''

\subsection{Karatsuba Modular Square ($-22.4$ Million Toffoli)}

The largest single Toffoli reduction in the project's history.

\textbf{The schoolbook problem}: Computing $\lambda^2$ for a 256-bit number using schoolbook
multiplication costs $O(n^2)$ multiplications.

\textbf{Karatsuba}: For a 256-bit number split as $\lambda = \text{hi} \cdot 2^{128} + \text{lo}$:
\[
  \lambda^2 = \text{hi}^2 \cdot 2^{256}
            + \left[(\text{lo}+\text{hi})^2 - \text{lo}^2 - \text{hi}^2\right] \cdot 2^{128}
            + \text{lo}^2
\]
Instead of computing 4 products, compute only 3 squares: $\text{lo}^2$, $\text{hi}^2$,
$(\text{lo}+\text{hi})^2$. Reconstruct the cross-term via additions (cheap Clifford operations).

\textbf{Why $-22.4$M Toffoli}: The squaring operation runs at EVERY GCD step (258 steps $\times$
2 passes). Any $O(n^2) \to O(n^{1.585})$ improvement compounds massively across all those calls.

\textbf{Density-neutral}: The Karatsuba identity is a pure algebraic identity, correct for all
integers.

\textbf{Recursive Karatsuba}: Tested on 128-bit sub-squares $\to$ net Toffoli LOSS (the overhead
of the additions outweighs the sub-square savings at that level). Dead end.

\subsection{NAF Recoding of Constants}

\textbf{Non-Adjacent Form (NAF)} uses signed digits $\{-1, 0, +1\}$ to represent constants with
fewer non-zero digits:
\[
  977 = 2^{10} - 2^{5} - 2^{4} + 1 \quad \text{(4 signed terms vs 6 binary terms)}
\]

In NAF, no two adjacent digits are both non-zero, and the number of non-zero digits is minimal.
For $977$: 4 terms instead of 6, saving $\approx33\%$ of the Solinas fold cost.

\textbf{Density-neutral}: The identity $977 = 1024 - 32 - 16 + 1$ is a mathematical fact, correct
for all inputs.

\subsection{Structural Dead Gate Skipping}
\label{sec:deadccx}

Some CCX gates provably never fire because their control condition is structurally always zero.

\textbf{Two flavors} (with very different correctness profiles):

\begin{enumerate}
  \item \textbf{Empirical dead-CCX} (REMOVED from SOTA): Sample $\approx9$ million inputs; find
    CCX gates whose target never flips; skip them. Problem: these gates might fire on inputs not
    in the sample. Removed in commit \texttt{6dafa07}.

  \item \textbf{Structural dead-CCX} (CURRENT SOTA): Prove mathematically that a CCX gate never
    fires on ANY input using circuit structure alone (known-constant carries, exact-remainder
    conditions, call-site bit patterns). Named predicates in \texttt{trailmix\_ludicrous}:
    \texttt{TLM\_FFG\_SKIP\_STRUCTURAL\_DEAD\_*},
    \texttt{TLM\_COMPARE\_SKIP\_STRUCTURAL\_DEAD\_*}, etc. Density-neutral.
\end{enumerate}

\textbf{Iterated (two-pass) dead-CCX} (\texttt{DROP\_DEAD\_ROBUST\_SECOND}, the empirical flavor ---
this is what carried the 1153 SOTA over the line before the structural pivot): dead-CCX deletion is
\textbf{iterable}. Drop a first list, then \textbf{re-screen the already-dropped op stream} --- the
first drop reshapes the stream, so a second screen finds gates that are newly (or only now
detectably) dead. Drop a second, smaller list and re-hunt the nonce with \emph{both} drops on. Each
pass is subset-monotone and the returns diminish, but they are real (the second pass was worth
$\sim$$-2{,}411$ avgT at 1153q). This is a distribution-exact / island-overfit lever
(\S\ref{sec:density}): the indices are absolute positions in one exact op stream, so any structural
change invalidates them and forces a full re-screen. Reusable mainly for squeezing a frozen final
artifact; the durable wins are the \emph{structural} levers above.

\subsection{Classical Constant Folding}

When one operand of an arithmetic operation is a \textbf{classical constant} (known at compile
time), the operation can often be implemented with 0 Toffoli gates.

\textbf{The mechanism}: A \texttt{BitId} register holds a constant pattern as compile-time
information. Operations with one classical operand:
\begin{itemize}
  \item \texttt{qubit XOR 0} = identity (no gate)
  \item \texttt{qubit XOR 1} = Pauli $X$ (free Clifford gate)
  \item \texttt{qubit + classical\_constant mod p} = a series of CNOT-like operations --- 0 Toffoli.
\end{itemize}

\textbf{In trailmix\_ludicrous}: $Q = (Q_x, Q_y)$ is classical, so \texttt{coord\_add3x(circ, x2,
ox, qubits)} (\textit{x2 $\mathrel{+}$= 3$\cdot Q_x$ mod p}) costs \textbf{0 Toffoli}.

\subsection{Converged-Tail Gate Elision}

In the last $K$ iterations of the 530-step GCD, the algorithm has converged and the apply-phase
\texttt{cswap} has a control that is deterministically 0 on all but rare inputs. Skip the cswap
for those last $K$ iterations.

\textbf{Cost}: Island-exact. Requires re-hunting the tail nonce after applying.

\subsection{GAP\_J2 Comparator Window Narrowing ($-1.33$M Toffoli)}

The swap decision comparator at each GCD step checks whether $|u| > |v|$. A full 256-bit
comparison would be expensive, but the GCD ensures that $u$ and $v$ typically differ somewhere
in their top bits.

\textbf{The schedule \texttt{GAP\_J2[i]}} (258-element table) sets the comparator window width
per GCD step. The comparator scans only the top $\texttt{GAP\_J2[i]}$ bits. A 22-line edit to the
table to shave $\approx1$ bit per step across 258 steps:
\[
  \text{Toffoli saved} = 1\ \text{bit} \times 516\ \text{comparator calls} \approx 1.33\text{M Toffoli}
\]
\textbf{Island-exact} --- requires nonce hunting.

\subsection{Ancilla-Free Majority}
\label{sec:majority}

Standard carry logic uses a 3-input majority gate with 3 Toffoli and an ancilla. The identity:
\[
  \text{maj}(a, k, c) = c \oplus ((a \oplus c) \wedge (k \oplus c))
\]
implements the same majority with \textbf{2 Toffoli and no ancilla} (XOR $=$ CNOT, AND $=$ CCX).
Peak-neutral, value-identical, pure count reduction wherever majority gates appear.

\subsection{Exact-Adder Recovery}

Counter-intuitive: Sometimes a truncated adder costs \textbf{more} Toffoli than the exact
full-carry adder, because the truncation requires an expensive correction to handle the rare
carry-escape cases.

\textbf{The rule}: If the correction overhead of a truncated adder exceeds its savings from
dropping bits, revert to the exact adder. Named: \texttt{DIALOG\_GCD\_APPLY\_FINAL\_TOPCLEAN=0}
$\to$ $\approx{-2{,}597}$ avg-Toffoli, value-exact.

\subsection{Skip Redundant Final Cleanup}
\label{sec:skipclean}

A folded/structured computation sometimes makes a final ``top-clean'' pass redundant --- the
structure ensures the cleanup bits are already zero by construction. Dropping the cleanup removes
its Toffoli.

\begin{tcolorbox}[noteboxstyle, title=Highest-risk move]
If the cleanup is not actually redundant, you get silent phase/classical dirt on some inputs.
Always pair with strong \cls/\pha/\anc{} validation before enabling.
\end{tcolorbox}

\subsection{Off-Peak Loosening (The Dual Move)}

When a phase falls \textbf{off} the qubit peak (another phase is now the binder), you can WIDEN
that phase's arithmetic to recover Toffoli for free:
\begin{lstlisting}
Phase A: was binder at 1170q -> locally reduced to 1162q
Phase B: now binder at 1170q (unchanged)
Action: widen Phase A's carry width back toward the old value
Effect: Phase A's Toffoli drops; Peak stays at 1170 (Phase B is binding)
\end{lstlisting}

\textbf{The rule}: Shrinking an off-peak phase is a pure-Toffoli play. Widening an off-peak phase
is a free Toffoli refund. Govern both on gates alone, not qubits.

\subsection{Constprop: Automated Static Gate Elimination}

The constraint-propagation pass (\texttt{CONSTPROP\_MAX\_ITERS}) analyzes the circuit statically
and eliminates:
\begin{itemize}
  \item \textbf{Provably-constant CCX gates}: control always 0 $\to$ remove; control always 1
    $\to$ replace with CNOT
  \item \textbf{Inverse-pair cancellation}: consecutive \texttt{CCX(a,b,c)} gates that cancel
    algebraically
  \item \textbf{Affine simplification}: sequences simplifying to an affine function, implementable
    with fewer gates
\end{itemize}

Setting \texttt{CONSTPROP\_MAX\_ITERS=256} (vs default 16) runs to a deeper fixpoint, finding
more opportunities: $-377$k Toffoli in one version.

\subsection{Quadratic-Cascade $\to$ Controlled-Increment (cinc)}
\label{sec:cinc}

\textbf{Replace an $O(t^2)$ gadget with an $O(t)$ one when the math allows.} A recurring pattern is
a ``carry-into-tail'' or prefix-propagation built as a \textbf{quadratic cascade of
multi-controlled-X gates} (\lstinline{mcx_clean_k} over a window $[nv, L)$ --- every position
controls on all positions below it, giving $O(t^2)$ MCX work).

When the cascade is computing a \emph{carry propagation} (the standard ``increment if all low bits
are 1'' pattern), it is exactly a \textbf{controlled increment}, and there is a cheaper exact
primitive for that: the \textbf{Khattar--Gidney controlled-increment} \texttt{cinc}
(arXiv:2407.17966), which uses only $\log^* n$ clean ancilla and $O(n)$ Toffoli instead of the
quadratic cascade.

\begin{lstlisting}
Old: clean-tail fold carry = mcx_clean_k prefix cascade over [nv, L)   -> O(t^2) MCX
New: clean-tail fold carry = cinc_khattar_gidney (one controlled +1)   -> O(t), log*-ancilla
\end{lstlisting}

\textbf{Value-exact and density-neutral}: an exact functional replacement (same truth table), so it
does not touch the island distribution --- it can be stacked freely without re-hunting a nonce.
Named \texttt{TLM\_FOLD\_TAIL\_CINC}; worth $\sim$$-5{,}175$ avgT alone. This is the \textbf{``safe''
Toffoli lever class}: unlike dead-CCX deletion (\S\ref{sec:deadccx}, distribution-exact), a cinc swap
is correct on all inputs --- prefer this class when choosing what to stack first.

The general principle: \textbf{whenever a gadget's cost is quadratic in a window width, ask whether
it is secretly a standard primitive (increment, compare, prefix-AND) with a known linear
construction.}

% ===========================================================================
\section{The Qubit $\leftrightarrow$ Toffoli Exchange Rate}
\label{sec:exchange}
% ===========================================================================

\subsection{The fundamental arithmetic}

At the current SOTA ($1152\text{q} \times 1{,}364{,}230\ \text{avgT} = 1{,}571{,}592{,}960\ \text{score}$):
\[
  \text{break-even rate} = \frac{\text{avgT}}{\text{peak\_q}} = \frac{1{,}364{,}230}{1152} \approx 1{,}184\ \text{Toffoli per qubit}
\]
A lever that saves 1 qubit at the cost of fewer than 1,184 Toffoli \textbf{improves the score}.
A lever that costs more than 1,184 Toffoli per qubit saved \textbf{worsens the score} --- even
though it reduced qubits.

This rate \textbf{changes with every optimization}:
\begin{itemize}
  \item After a large Toffoli reduction: the rate drops (qubits become more valuable). Previously-marginal qubit cuts are now worth trying.
  \item After a large qubit reduction: the rate rises. Previously-marginal Toffoli cuts might now dominate.
\end{itemize}

\subsection{The product-race caveat}

A qubit drop that clears break-even can STILL lose the product race if the Toffoli-grinding
track moves faster in parallel.

\textbf{Concrete example} (June 2026):
A 1157q clamp cost $\approx1{,}127$ Toffoli/qubit --- below break-even at landing. But the 1159q
Toffoli-grind reached $1{,}378{,}242$ avgT. Score comparison:
\[
  1159 \times 1{,}378{,}242 = 1{,}597{,}506{,}378 < 1157 \times 1{,}380{,}890 = 1{,}598{,}290{,}330
\]
The 2-qubit savings was overtaken by the parallel Toffoli-cutting track.

\textbf{Lesson}: Always run both tracks and compare products. Never minimize qubit count in isolation.

\subsection{The shelved-lever rule}

A qubit lever that \textbf{failed break-even on one Toffoli base} can become favorable on a
cheaper Toffoli base.

\textbf{Example}: The dynamic headroom clamp (\texttt{TLM\_TARGET\_Q}) cost $\approx1{,}514$
Toffoli/qubit on the schoolbook-square base. After Karatsuba reduced the square cost, the same
clamp cost only $\approx20$ Toffoli/qubit on the new base --- well below break-even. It won.

\textbf{Rule}: After any major arithmetic win (Karatsuba, algebraic fusion, structural rewrite),
re-test \textbf{all previously-shelved qubit levers} at the new exchange rate.

\subsection{Different peak layers have different exchange rates}

At any peak, there are typically two layers:

\begin{enumerate}
  \item \textbf{Persistent floor}: long-lived state held across the whole peak (transcript,
    passenger values, GCD state). Often \textbf{cheap} to reduce --- a denser codec or transcript
    streaming costs few gates to save many qubits.
  \item \textbf{Transient working registers}: the current step's carry bits, arithmetic scratch.
    Often \textbf{expensive} to reduce --- freeing them requires recompute/dirty-borrow, adding
    many Toffoli.
\end{enumerate}

Counter-intuitively: \textbf{compress the transcript before freeing the scratch}. The transcript
looks ``irreducible'' (structured data) but is often the cheap lever. The scratch looks
``wasteful'' but is often the expensive lever.

% ===========================================================================
\section{Density-Neutral vs Island-Exact}
\label{sec:density}
% ===========================================================================

\subsection{The fundamental distinction}

\begin{itemize}
  \item \textbf{Density-neutral (value-exact)}: The circuit produces the same output on ALL
    inputs. The optimization is grounded in mathematical identity or circuit structure.
  \item \textbf{Island-exact (distribution-specific)}: The circuit is correct only on the
    specific 9024 test inputs. On other inputs, it may give wrong answers.
\end{itemize}

When you apply an island-exact optimization, you need to find a new \textbf{nonce} --- a 48-byte
seed for the test input generator --- such that all 9024 generated inputs happen to be in the
safe zone. Stacking multiple island-exact optimizations makes the safe zone smaller.

\subsection{The 6dafa07 pivot: empirical $\to$ structural dead-CCX}

Before the current SOTA, the circuit used \textbf{empirical dead-CCX elimination}: sample 9M
inputs, find CCX gates that never fire, skip them. Island-exact.

The \texttt{6dafa07} commit replaced this with \textbf{structural dead-CCX skipping}: prove from
circuit structure which CCX gates can never fire on \emph{any} input. These predicates are keyed
by call-site and bit-index, not sampled data. Density-neutral.

\subsection{The correctness check triple: \cls/\pha/\anc}

Every validation run checks three error types:
\begin{itemize}
  \item \cls: classical mismatch (the output value is wrong on some input)
  \item \pha: phase error (the output is correct but entangled with extra phase --- kills interference)
  \item \anc: ancilla mismatch (a qubit failed to uncompute back to \ket{0})
\end{itemize}

A correct circuit shows \texttt{0/0/0}. Measurement-based uncompute fails in \pha{} (which value
checks miss). Truncation failures typically show in \cls. Missed uncompute shows in \anc.

% ===========================================================================
\section{Pebbling Theory}
\label{sec:pebbling}
% ===========================================================================

\subsection{The reversible pebble game}

The \textbf{pebble game} models space-time tradeoffs in reversible computation. The computation is
a directed graph (nodes $=$ values, edges $=$ dependencies). A ``pebble'' $=$ a qubit holding that
value. Peak qubits $=$ maximum pebbles simultaneously.

\subsection{Bennett pebbling: classical reversibility is expensive}

Bennett (1989) showed that reversibly simulating a $T$-step computation using $S$ extra pebbles
costs:
\[
  \text{Time} = \varepsilon \cdot 2^{1/\varepsilon} \cdot T^{1+\varepsilon} / S^\varepsilon \quad \text{(for any $\varepsilon > 0$)}
\]
The hidden constant $c(\varepsilon) = \varepsilon \cdot 2^{1/\varepsilon}$ grows \textbf{exponentially}
as $\varepsilon \to 0$.

\textbf{Concrete example}: Halving the GCD transcript ($S$ halved, $\varepsilon = 1$):
\[
  \text{Time blowup} \approx 1 \times 2 \times 530^2 / 370 \approx 1{,}518\ \text{extra GCD passes}
\]
At $\approx2{,}500$ Toffoli per pass: 3.8M extra Toffoli for $\approx370$q savings.
Break-even needs: $3.8\text{M} / 1{,}184 \approx 3{,}200$ qubits. Actual saving: 370. Not worth it.

\subsection{Quantum (spooky) pebbling: exponentially better}

\textbf{Gidney (2019)} introduced \textbf{spooky pebbling}: also allow removing a pebble by
measuring it in the $X$-basis, recording the measurement outcome as a classical bit. When the
value is needed again, recompute it and apply a classically-conditioned phase correction.

\textbf{Kornerup, Sadun, Soloveichik (2021)} proved tight bounds:
\begin{align*}
  \text{Quantum:}   &\quad T_{\text{quantum}} = O(T/\varepsilon)           \quad [\text{polynomial constant}] \\
  \text{Classical:} &\quad T_{\text{classical}} = O(2^{1/\varepsilon} \cdot T) \quad [\text{exponential constant}]
\end{align*}
The quantum advantage is exponential. This is the theoretical backbone of why MBU and
measurement-based techniques dominate.

\subsection{Parallel spooky pebbling (Kahanamoku-Meyer et al.\ 2025)}

For a length-$\ell$ sequential computation chain:
\[
  \text{qubits needed} = 2.47 \times \log(\ell) \quad \text{at depth } 2\ell
\]

For our 530-step GCD: $2.47 \times \log(530) \approx 23$ qubits. Our current GCD state machine
uses $\approx22$ qubits --- already near the theoretical minimum.

% ===========================================================================
\section{The SOTA Circuit: How trailmix\_ludicrous Works}
\label{sec:trailmix}
% ===========================================================================

\subsection{The decisive architectural choice: Q is classical}

The second point $Q = (Q_x, Q_y)$ is kept as classical \texttt{BitId} values, not quantum
registers. Materialized into a transient quantum temp \emph{only} at off-peak coordinate steps,
then freed. This eliminates 512 qubits from the peak.

\subsection{Two GCD passes sharing one tape}

The affine point addition needs:
\begin{enumerate}
  \item $\lambda = dy \cdot dx^{-1} \bmod p$ --- inverse GCD (drive $(p, dx) \to (1, 0)$)
  \item $\text{new\_y\_component} = \lambda \cdot (P_x - P_x') \bmod p$ --- forward GCD
\end{enumerate}

Both share \textbf{one} modular-inverse engine and \textbf{one} dialog tape (transcript). This
avoids having two separate inversion circuits, saving both qubits (one shared tape) and Toffoli
(shared codec overhead).

\subsection{The base-5 codec: 772 $\to$ 603 qubits live}

At each of 257 divstep steps, the algorithm records $(\texttt{subtracted}, \texttt{swap}, s_2)$ ---
3 bits raw. But only 5 of 8 possible 3-bit patterns are reachable. Triple windows have $5^3 = 125$
reachable states, encoded in 7 bits instead of 9:
\begin{lstlisting}
Layout: 1 Step0 (2b) + 85 Triples (7b each) + 1 Pair tail (5b) + 1 t1 flag (1b) = 603 qubits
\end{lstlisting}

\subsection{The vented-adder zoo}

Every adder in \texttt{trailmix\_ludicrous} vents its carry-uncompute via MBU:
\begin{lstlisting}[language=Rust]
hmr(carry, bit);       // X-basis measurement, 0 Toffoli
cz_if_bit(a, b, bit);  // phase correction conditioned on measurement
\end{lstlisting}

The adder family dispatches per-call by baked schedule tables. Each call's vent budget: exactly
\texttt{TLM\_TARGET\_Q} $-$ \texttt{active\_qubits}, spending every spare qubit below the ceiling
on Toffoli savings.

\subsection{The deliberate island-exact components}

trailmix\_ludicrous includes calibrated island-exact approximations:
\begin{itemize}
  \item \textbf{Low-54-bit +f fold} (\texttt{LSBS = 54}): reduction folds touch only low 54 bits;
    carry beyond bit 53 is dropped with probability $\approx 2^{-21}$ per fire.
  \item \textbf{Narrow-top-window comparators} (\texttt{MSBS = PAD = 21}): swap predicates scan
    only the top 21 bits. Mis-decision probability $\approx 2^{-21}$ per call.
\end{itemize}

\subsection{Where 1152q comes from}

At the SOTA (\texttt{d44cad3}):
\begin{itemize}
  \item \textbf{603q}: compressed base-5 GCD transcript (live during both GCD passes)
  \item \textbf{512q}: two 256-bit coordinate registers (quantum $P_x$, $P_y$)
  \item \textbf{$\approx37$q}: transient carry/vent ancilla at the peak-binding fold operations
\end{itemize}

% ===========================================================================
\section{The 1211$\to$1152 Historical Journey}
\label{sec:history}
% ===========================================================================

\begin{center}
\begin{tabular}{p{3.8cm}p{4.5cm}p{5.2cm}}
\toprule
\textbf{Transition} & \textbf{Technique} & \textbf{What was released} \\
\midrule
1211$\to$1193 ($-18$q) & Live-range hole (\S\ref{sec:liverange}) & GCD-apply scratch block freed during shift sub-phase; reacquired after \\
1193$\to$1192 ($-1$q) & Truncation (\S6.D) & Square segment one notch tighter \\
1192$\to$1185 ($-7$q) & Transcript codec (\S\ref{sec:transcript}) & Per-step fold carry scheduling + exact 11-bit head codec \\
1185$\to$1170 ($-15$q) & Plateau decomp (\S6.I) + 5 techniques & Tail codec, streaming apply, square rebalance, 5 co-binders \\
1170$\to$1169 & MBU vent on fold carry (\S\ref{sec:mbu}/8.B) & \texttt{hmr+cz\_if} on the $-f$ fold peak owner: 0 Toffoli to vent it \\
1167$\to$1166 & Step-0 swap\_flag elimination & \texttt{swap\_flag} becomes \texttt{Option<QubitId>}, lazy-alloc; freed one tape lane \\
1166$\to$1165 & Odd-u0 parking (\S\ref{sec:liverange}) & \texttt{u[0]} is provably 1; park+loan the lane, restore in 2 gates \\
1164$\to$1163 & Source-as-carry suffix & Paid drop: surrendered apply-phase vents to free one carry lane \\
Toffoli: $-22.4$M & Karatsuba square (\S8.F) & 3 sub-squares instead of 4 \\
Toffoli: $-1.33$M & GAP\_J2 narrowing (\S8.K) & 22-line per-step comparator table edit \\
Toffoli: $-377$k & Constprop deeper (\S8.P) & \texttt{CONSTPROP\_MAX\_ITERS} $16\to256$ \\
1163$\to$1159$\to$1156 & Headroom clamp returns (\S6.J + shelved-lever \S\ref{sec:exchange}) & Same clamp that lost early now wins on the cheap Karatsuba$+$dead-CCX base ($\sim$527 T/qubit) \\
1156$\to$1153 & Clamp $+$ zero-cin fold $+$ FFG cap (\S6.J), paid by cinc (\S\ref{sec:cinc}) $+$ iterated dead-CCX (\S\ref{sec:deadccx}) & Qubit drop alone cost $+11{,}369$ avgT (rejected); Toffoli levers stacked back to win \\
\textbf{1153$\to$1152} & Free-and-recompute cy0 (\S6.E) & \texttt{cy0 = ctrl \& !a0}; loan lane for 40 binding folds \\
Structural & 6dafa07 pivot (\S8.H/\S\ref{sec:density}) & Replaced empirical dead-CCX with $\approx15$ structural predicates \\
\bottomrule
\end{tabular}
\end{center}

\textbf{Key meta-lessons}:
\begin{enumerate}
  \item \textbf{Every ``floor'' fell to a better encoding or a better question.} Not always a new
    algorithm --- often a denser codec or the question ``whose value is a cheap function of live
    qubits?''

  \item \textbf{The shelved-lever rule is real.} The 1157q clamp tried early and lost; won after Karatsuba.

  \item \textbf{Plateau mode is the endgame.} The final qubit drops each required coordinated
    packages of 4--9 independent local reductions applied simultaneously.

  \item \textbf{Structural correctness wins over island-exact tricks.} The 6dafa07 pivot to
    structural dead-CCX eliminated the empirical approach entirely --- correct for all inputs,
    not just the 9024.
\end{enumerate}

\subsection{The low-qubit competition track (a separate objective --- \S\ref{sec:twocomps})}

Running \emph{alongside} the Q$\times$T product-SOTA journey above is the separate \textbf{low-qubit
competition} line --- the \textbf{Shrunken-PZ / Proos--Zalka divstep} family --- whose objective is
to minimize peak qubits with Toffoli unconstrained. This is the frontier where EEA register sharing
(\S\ref{sec:regshare}) shines:

\begin{center}
\begin{tabular}{p{2.6cm}p{6.0cm}p{4.4cm}}
\toprule
\textbf{Qubit record} & \textbf{Technique} & \textbf{Notes} \\
\midrule
1050q & trailmix shrunken-PZ embedded op stream & early low-qubit route \\
1019q & source-level dynamic-width PZ port & \\
988$\to$956q & dirty-borrow / gate-hosting lever stack & \\
952$\to$948q & further lane borrowing + thin schedule & lowest \emph{fully-validated} witness \\
\textbf{851q} (\texttt{e7dd3de}, 6/23) & \textbf{EEA register sharing (\S\ref{sec:regshare}), Luo 2026} & \textbf{$-97$q; the breakthrough} \\
829q (\texttt{1dd61ca}, 6/24) & more borrowing/relocation on a frozen envelope & current low-qubit frontier \\
\bottomrule
\end{tabular}
\end{center}

These are records on the \textbf{pure-qubit objective}, where the $\sim$400--500M Toffoli they spend
simply does not enter the scoring. The 851q and 829q figures come from an analysis-oracle accounting
(\S\ref{sec:regshare}) --- read them as the leading \textbf{lower-bound witnesses} for the qubit
floor; the 948q route is the lowest fully-validated circuit on the track. This is a \emph{different
game} from the Q$\times$T product SOTA (1152q $\times$ 1.36M, value-exact): more qubits but
$\sim$340$\times$ fewer Toffoli. The two minima are reached by different constructions, and which one
is ``best'' depends entirely on which competition you are entering --- see \S\ref{sec:regshare} and
\S\ref{sec:exchange}.

% ===========================================================================
\section{Open Frontiers}
\label{sec:frontiers}
% ===========================================================================

\subsection{Density-neutral Toffoli cuts not yet applied}

\begin{enumerate}
  \item \textbf{Gidney 2025 streaming vented adder} (arXiv:2507.23079): Uses 2 clean $+$ $(n-2)$
    dirty ancilla for $\approx3n$ Toffoli. Partially implemented in \texttt{venting.rs} but
    \textbf{leaks phase} --- dirty ancilla phase corrections are incomplete. Expected $\approx n/4$
    Toffoli savings per adder once fixed.

  \item \textbf{Apply-swap structural dead-CCX}: The GCD apply-step \texttt{cswap} block
    ($\approx256$ cswap per iter $\times$ 258 iters $\times$ 2 passes) has NO per-bit
    structural-dead skip. Porting the structural-dead analysis to these operations is a
    high-priority open lever.

  \item \textbf{Extended Cuccaro structural dead-carry}: Currently checked for call indices 13--37
    only. Extending to ALL Cuccaro adder call sites would save proportionally more.

  \item \textbf{Apply-swap involutory pair cancellation}: The forward GCD-swap and the immediately
    following GCD-cswap on the same $(u/v)$ lanes may cancel algebraically (swap twice $=$ identity).
    Not yet checked structurally.
\end{enumerate}

\subsection{Density-neutral qubit cuts not yet applied}

\begin{enumerate}
  \item \textbf{The co-bind plateau at 1152q}: Multiple independent phases all tie at 1152q.
    Each needs a local reduction before the global peak drops.

  \item \textbf{FFG cy0-style recompute on other binders}: The cy0 trick (\S6.E) worked because
    \texttt{cy0} was a cheap function of live qubits. Other binder qubits may have similar
    cheap-function properties not yet identified.

  \item \textbf{SumHiLo/shifted vents} (\texttt{TLM\_SQUARE\_SUMHILO\_VENT},
    \texttt{TLM\_SQUARE\_VENT\_SHIFTED}): Exist in the codebase as default-OFF options. Off-peak
    phases could enable these for free Toffoli savings.
\end{enumerate}

% ===========================================================================
\section{Quick-Reference Summary Tables}
\label{sec:tables}
% ===========================================================================

\subsection{Qubit reduction techniques}

\begin{center}
\small
\begin{tabular}{p{4.8cm}p{3cm}p{2.8cm}p{2.5cm}}
\toprule
\textbf{Technique} & \textbf{Qubits saved} & \textbf{Toffoli cost} & \textbf{Density-neutral} \\
\midrule
Live-range hole (\S6.A) & $k$ qubits & $+2 \times C_V$ & Yes \\
Passenger relocation (\S6.B) & $k$ qubits & 0 & Yes \\
Range-bound guard removal (\S6.C) & 1 per register & 0 & Yes \\
Truncation (\S6.D) & varies & 0 & Partial \\
Free-and-recompute cheap value (\S6.E) & 1 & $\approx0$ (CNOTs) & Yes \\
Dynamic-width registers (\S6.F) & varies & small & Partial \\
Known-constant teardown (\S6.G) & $k$ & 0 (CNOT $=$ Clifford) & Yes \\
Transcript compression (\S6.H) & varies & decode overhead & Yes \\
Plateau decomposition (\S6.I) & 1 (global) & combined cost & Yes \\
Dynamic headroom clamp (\S6.J) & 1 per ceiling drop & small & Yes \\
Lazy carry liveness (\S6.K) & 1/chunk & 0 & Yes \\
In-place decode (\S6.L) & varies & codec overhead & Yes \\
Borrowed-carry fusion (\S6.M) & $k$ & 0 & Yes \\
HMR ghosting (\S6.N) & 256 & $+1$ modular multiply & Yes \\
EEA register sharing (\S6.O) & $\sim$97--119 (to 851$\to$829q) & high (irrelevant to low-q objective) & Yes --- low-qubit competition \\
Gate-hosting (\S6.P) & 1 per hosted gate & 0 (clean) / $+1$ Tof (dirty) & Yes if host-zero \emph{proven}; island-exact if \emph{sampled} \\
\bottomrule
\end{tabular}
\end{center}

\subsection{Toffoli reduction techniques}

\begin{center}
\small
\begin{tabular}{p{4.8cm}p{3cm}p{2.5cm}p{2.5cm}}
\toprule
\textbf{Technique} & \textbf{Toffoli saved} & \textbf{Qubit cost} & \textbf{Density-neutral} \\
\midrule
MBU / carry vent (\S8.A/B) & 1 per vent & $+1$ temporary & Yes \\
Conditional replay (\S8.C) & $(1-f)\times\text{block}$ & 0 & Yes \\
Algebraic fusion (\S8.D) & varies & 0 & Yes \\
Witness-first dataflow (\S8.E) & $\approx13$M / instance & 0 & Yes \\
Karatsuba square (\S8.F) & $\approx22$M (one change) & 0 & Yes \\
NAF recoding (\S8.G) & $\approx5$--10\% of fold cost & 0 & Yes \\
Structural dead-gate skip (\S8.H) & varies & 0 & Yes \\
Classical constant ops (\S8.I) & significant & 0 & Yes \\
Converged-tail elision (\S8.J) & hundreds/iter $\times K$ & 0 & No \\
GAP\_J2 narrowing (\S8.K) & $\approx1.33$M / 22-line edit & 0 & No \\
Ancilla-free majority (\S8.L) & 1 per majority & 0 & Yes \\
Exact-adder recovery (\S8.M) & $\approx2{,}600$/call & 0 & Yes \\
Redundant cleanup skip (\S8.N) & varies & 0 & Risky \\
Off-peak loosening (\S8.O) & free refund & 0 & Yes \\
Constprop deeper (\S8.P) & $\approx377$k & 0 & Yes \\
Cinc cascade replacement (\S8.Q) & $\approx5{,}175$ ($O(t^2)\to O(t)$) & 0 & Yes \\
Iterated 2-pass dead-CCX (\S8.H) & $\approx2{,}400$/extra pass & 0 & No \\
\bottomrule
\end{tabular}
\end{center}

\subsection{Key theoretical results}

\begin{center}
\small
\begin{tabular}{p{4.5cm}p{4.5cm}p{4.5cm}}
\toprule
\textbf{Result} & \textbf{Source} & \textbf{Insight} \\
\midrule
MBU: uncompute AND for 0 $T$-gates & Jones 2013 (Phys Rev A); Gidney 2018 (arXiv:1709.06648) & $X$-measurement $+$ CZ $=$ free uncompute \\
1 Toffoli $=$ exactly 7 $T$-gates & Gosset et al.\ 2013 (arXiv:1308.4134) & Tight lower bound \\
Bennett classical pebbling & Bennett 1989 (SIAM J.\ Comput.\ 18:4) & Classical: exponential constant in space savings \\
Quantum spooky pebbling & Kornerup, Sadun, Soloveichik 2021 (arXiv:2110.08973) & Quantum: polynomial constant (exponential advantage) \\
Parallel spooky pebbling & Kahanamoku-Meyer et al.\ 2025 (arXiv:2510.08432) & $2.47 \times \log(\ell)$ qubits for $\ell$-step chain \\
Conditionally-clean ancilla & Khattar \& Gidney 2024 (arXiv:2407.17966) & $3n$ Toffoli MCX with $\log^* n$ clean ancilla \\
secp256k1 ECDLP frontier & Babbush, Gidney et al.\ 2026 (arXiv:2603.28846) & $\leq1200$q / $\leq90$M Toffoli for 256-bit ECDLP \\
Karatsuba algorithm & Karatsuba \& Ofman 1962 & $O(n^{1.585})$ vs $O(n^2)$ for multiplication \\
EEA register sharing (origin) & Proos \& Zalka 2003 (arXiv:quant-ph/0301141) & Inversion dominates point-add space; operand$+$cofactor packing \\
Compact exact reversible inversion & Luo et al.\ 2026 (arXiv:2604.02311) & $3n + 4\lfloor\log_2 n\rfloor$ qubits (1333q, $n{=}256$) via bit-length-tracked sharing \\
\bottomrule
\end{tabular}
\end{center}

\vfill
\noindent\rule{\textwidth}{0.4pt}
\vspace{2pt}

\noindent\textit{See also} (all under
\path{skills/ecdsafail_skills/ecdsafail-circuit-optimization/references/}):
\begin{itemize}\small
  \item \texttt{density\_neutral\_tradeoffs.md} --- detailed technique catalog with exchange rate math
  \item \texttt{gidney-techniques.md} --- Gidney's full lever catalog with blog/paper pointers
  \item \texttt{external-literature-2000-2026.md} --- broader academic literature map
  \item \texttt{REPORT\_1168\_wall\_revamp.md} --- the trailmix\_ludicrous introduction
  \item \texttt{frontier-1211-to-1170.md} --- detailed 1211$\to$1170 step-by-step record
\end{itemize}

\end{document}
